\documentclass{amsart}
% For dealing with references we use the comment environment
\usepackage{verbatim}
\newenvironment{reference}{\comment}{\endcomment}
%\newenvironment{reference}{}{}
\newenvironment{slogan}{\comment}{\endcomment}
\newenvironment{history}{\comment}{\endcomment}

% For commutative diagrams we use Xy-pic
\usepackage[all]{xy}

% We use 2cell for 2-commutative diagrams.
\xyoption{2cell}
\UseAllTwocells

% We use multicol for the list of chapters between chapters
\usepackage{multicol}

% This is generally recommended for better output
\usepackage{lmodern}
\usepackage[T1]{fontenc}

% For cross-file-references

% Package for hypertext links:
\usepackage{hyperref}

% For any local file, say "hello.tex" you want to link to please

% Theorem environments.
%
\theoremstyle{plain}
\newtheorem{theorem}[subsection]{Theorem}
\newtheorem{proposition}[subsection]{Proposition}
\newtheorem{lemma}[subsection]{Lemma}

\theoremstyle{definition}
\newtheorem{definition}[subsection]{Definition}
\newtheorem{example}[subsection]{Example}
\newtheorem{exercise}[subsection]{Exercise}
\newtheorem{situation}[subsection]{Situation}

\theoremstyle{remark}
\newtheorem{remark}[subsection]{Remark}
\newtheorem{remarks}[subsection]{Remarks}

\numberwithin{equation}{subsection}

% Macros
%
\def\lim{\mathop{\mathrm{lim}}\nolimits}
\def\colim{\mathop{\mathrm{colim}}\nolimits}
\def\Spec{\mathop{\mathrm{Spec}}}
\def\Hom{\mathop{\mathrm{Hom}}\nolimits}
\def\Ext{\mathop{\mathrm{Ext}}\nolimits}
\def\SheafHom{\mathop{\mathcal{H}\!\mathit{om}}\nolimits}
\def\SheafExt{\mathop{\mathcal{E}\!\mathit{xt}}\nolimits}
\def\Sch{\mathit{Sch}}
\def\Mor{\mathop{\mathrm{Mor}}\nolimits}
\def\Ob{\mathop{\mathrm{Ob}}\nolimits}
\def\Sh{\mathop{\mathit{Sh}}\nolimits}
\def\NL{\mathop{N\!L}\nolimits}
\def\CH{\mathop{\mathrm{CH}}\nolimits}
\def\proetale{{pro\text{-}\acute{e}tale}}
\def\etale{{\acute{e}tale}}
\def\QCoh{\mathit{QCoh}}
\def\Ker{\mathop{\mathrm{Ker}}}
\def\Im{\mathop{\mathrm{Im}}}
\def\Coker{\mathop{\mathrm{Coker}}}
\def\Coim{\mathop{\mathrm{Coim}}}

% Boxtimes
%
\DeclareMathSymbol{\boxtimes}{\mathbin}{AMSa}{"02}

%
% Macros for moduli stacks/spaces
%
\def\QCohstack{\mathcal{QC}\!\mathit{oh}}
\def\Cohstack{\mathcal{C}\!\mathit{oh}}
\def\Spacesstack{\mathcal{S}\!\mathit{paces}}
\def\Quotfunctor{\mathrm{Quot}}
\def\Hilbfunctor{\mathrm{Hilb}}
\def\Curvesstack{\mathcal{C}\!\mathit{urves}}
\def\Polarizedstack{\mathcal{P}\!\mathit{olarized}}
\def\Complexesstack{\mathcal{C}\!\mathit{omplexes}}
% \Pic is the operator that assigns to X its picard group, usage \Pic(X)
% \Picardstack_{X/B} denotes the Picard stack of X over B
% \Picardfunctor_{X/B} denotes the Picard functor of X over B
\def\Pic{\mathop{\mathrm{Pic}}\nolimits}
\def\Picardstack{\mathcal{P}\!\mathit{ic}}
\def\Picardfunctor{\mathrm{Pic}}
\def\Deformationcategory{\mathcal{D}\!\mathit{ef}}

\setlength{\emergencystretch}{3em}
\begin{document}
\title[Stacks Project, Tag 0A3M]{278 --- Finite-dimensional top coherent cohomology of separated finite-type schemes over a field}
\maketitle
\noindent Copyright (C) 2005--2025 Johan de Jong. Stacks Project authors. Source: \href{https://stacks.math.columbia.edu/tag/0A3M}{Tag 0A3M}.

\noindent Editorial difficulty note (not author text). Difficulty H2: The statement does not assume properness, so standard proper cohomology finiteness cannot be applied directly. The proof combines induction on dimension and coherent devissage with Chow modification, relative ampleness and a lower-dimensional cokernel, then obtains a surjection from cohomology on a projective compactification..

\noindent Editorial provenance note (not author text). History: excerpt from revision \texttt{a0444\allowbreak 6e57e\allowbreak c1fbc\allowbreak 252a8\allowbreak 71afc\allowbreak ec775\allowbreak 2fb28\allowbreak 07b14}, etale-cohomology.tex, lines 9709--9835. Reference presentation was adapted; original-author context and core support blocks were retained or added. The mathematical wording is unchanged. Licensed under GNU FDL 1.2 or later; no invariant sections or cover texts. See ../licenses/COPYING. Context blocks reproduce author definitions and supporting results; they are not additional counted problems.

\begin{lemma}
\label{lemma-top-cohomology-coherent}
Let $X$ be a separated scheme of finite type over a field $k$.
Let $\mathcal{F}$ be a coherent sheaf of $\mathcal{O}_X$-modules.
Then $\dim_k H^d(X, \mathcal{F}) < \infty$ where $d = \dim(X)$.
\end{lemma}

\begin{proof}
We will prove this by induction on $d$. The case $d = 0$ holds because
in that case $X$ is the spectrum of a finite dimensional $k$-algebra $A$
(Varieties, Lemma \href{https://stacks.math.columbia.edu/tag/06LH}{lemma algebraic scheme dim 0 (Tag 06LH)})
and every coherent sheaf $\mathcal{F}$ corresponds to a finite $A$-module
$M = H^0(X, \mathcal{F})$ which has $\dim_k M < \infty$.

\medskip\noindent
Assume $d > 0$ and the result has been shown for separated schemes
of finite type of dimension $< d$. The scheme $X$ is Noetherian. Consider
the property $\mathcal{P}$ of coherent sheaves on $X$ defined by the rule
$$
\mathcal{P}(\mathcal{F}) \Leftrightarrow
\dim_k H^d(X, \mathcal{F}) < \infty
$$
We are going to use the result of
Cohomology of Schemes, Lemma \href{https://stacks.math.columbia.edu/tag/01YG}{lemma property initial (Tag 01YG)}
to prove that $\mathcal{P}$ holds for every coherent sheaf on $X$.

\medskip\noindent
Let
$$
0 \to \mathcal{F}_1 \to \mathcal{F} \to \mathcal{F}_2 \to 0
$$
be a short exact sequence of coherent sheaves on $X$.
Consider the long exact sequence of cohomology
$$
H^d(X, \mathcal{F}_1) \to
H^d(X, \mathcal{F}) \to
H^d(X, \mathcal{F}_2)
$$
Thus if $\mathcal{P}$ holds for $\mathcal{F}_1$ and $\mathcal{F}_2$,
then it holds for $\mathcal{F}$.

\medskip\noindent
Let $Z \subset X$ be an integral closed subscheme. Let $\mathcal{I}$
be a coherent sheaf of ideals on $Z$. To finish the proof we have to show
that $H^d(X, i_*\mathcal{I}) = H^d(Z, \mathcal{I})$ is finite dimensional.
If $\dim(Z) < d$, then the result holds because the cohomology group
will be zero (Cohomology, Proposition
\href{https://stacks.math.columbia.edu/tag/02UZ}{proposition vanishing Noetherian (Tag 02UZ)}).
In this way we reduce to the situation discussed in the following paragraph.

\medskip\noindent
Assume $X$ is a variety of dimension $d$ and 
$\mathcal{F} = \mathcal{I}$ is a coherent ideal sheaf. In this
case we have a short exact sequence
$$
0 \to \mathcal{I} \to \mathcal{O}_X \to i_*\mathcal{O}_Z \to 0
$$
where $i : Z \to X$ is the closed subscheme defined by $\mathcal{I}$.
By induction hypothesis we see that
$H^{d - 1}(Z, \mathcal{O}_Z) = H^{d - 1}(X, i_*\mathcal{O}_Z)$ is
finite dimensional. Thus we see that it suffices to prove the result
for the structure sheaf.

\medskip\noindent
We can apply Chow's lemma
(Cohomology of Schemes, Lemma \href{https://stacks.math.columbia.edu/tag/0200}{lemma chow Noetherian (Tag 0200)})
to the morphism $X \to \Spec(k)$. Thus we get a diagram
$$
\xymatrix{
X \ar[rd]_g & X' \ar[d]^-{g'} \ar[l]^\pi \ar[r]_i & \mathbf{P}^n_k \ar[dl] \\
& \Spec(k) &
}
$$
as in the statement of Chow's lemma. Also, let $U \subset X$ be
the dense open subscheme such that $\pi^{-1}(U) \to U$ is an isomorphism.
We may assume $X'$ is a variety as well, see
Cohomology of Schemes, Remark \href{https://stacks.math.columbia.edu/tag/0201}{remark chow Noetherian (Tag 0201)}.
The morphism $i' = (i, \pi) : X' \to \mathbf{P}^n_X$ is
a closed immersion (loc. cit.). Hence
$$
\mathcal{L} = i^*\mathcal{O}_{\mathbf{P}^n_k}(1) \cong
(i')^*\mathcal{O}_{\mathbf{P}^n_X}(1)
$$
is $\pi$-relatively ample (for example by
Morphisms, Lemma \href{https://stacks.math.columbia.edu/tag/02NR}{lemma characterize ample on finite type (Tag 02NR)}).
Hence by Cohomology of Schemes, Lemma \href{https://stacks.math.columbia.edu/tag/02O1}{lemma kill by twisting (Tag 02O1)}
there exists an $n \geq 0$ such that
$R^p\pi_*\mathcal{L}^{\otimes n} = 0$ for all $p > 0$.
Set $\mathcal{G} = \pi_*\mathcal{L}^{\otimes n}$.
Choose any nonzero global section $s$ of $\mathcal{L}^{\otimes n}$.
Since $\mathcal{G} = \pi_*\mathcal{L}^{\otimes n}$, the section $s$
corresponds to section of $\mathcal{G}$, i.e., a map
$\mathcal{O}_X \to \mathcal{G}$.
Since $s|_U \not = 0$ as $X'$ is a variety and $\mathcal{L}$
invertible, we see that $\mathcal{O}_X|_U \to \mathcal{G}|_U$
is nonzero. As $\mathcal{G}|_U = \mathcal{L}^{\otimes n}|_{\pi^{-1}(U)}$
is invertible we conclude that we have a short exact sequence
$$
0 \to \mathcal{O}_X \to \mathcal{G} \to \mathcal{Q} \to 0
$$
where $\mathcal{Q}$ is coherent and supported on a proper
closed subscheme of $X$. Arguing as before using our induction
hypothesis, we see that it
suffices to prove $\dim H^d(X, \mathcal{G}) < \infty$.

\medskip\noindent
By the Leray spectral sequence
(Cohomology, Lemma \href{https://stacks.math.columbia.edu/tag/01F4}{lemma apply Leray (Tag 01F4)})
we see that $H^d(X, \mathcal{G}) = H^d(X', \mathcal{L}^{\otimes n})$.
Let $\overline{X}' \subset \mathbf{P}^n_k$ be the closure
of $X'$. Then $\overline{X}'$ is a projective variety of dimension $d$
over $k$ and $X' \subset \overline{X}'$ is a dense open.
The invertible sheaf $\mathcal{L}^{\otimes n}$ is the restriction of
$\mathcal{O}_{\overline{X}'}(n)$ to $X'$. By
Cohomology, Proposition
\href{https://stacks.math.columbia.edu/tag/0A3G}{proposition cohomological dimension spectral (Tag 0A3G)}
the map
$$
H^d(\overline{X}', \mathcal{O}_{\overline{X}'}(n))
\longrightarrow
H^d(X', \mathcal{L}^{\otimes n})
$$
is surjective. Since the cohomology group on the left has
finite dimension by
Cohomology of Schemes, Lemma \href{https://stacks.math.columbia.edu/tag/01YS}{lemma coherent projective (Tag 01YS)}
the proof is complete.
\end{proof}
\end{document}
